
\scanpage{57}
\mspage{13}
On va peut-être se borner par la suite à des opérations de gommage de repères, ou de surmarquage, de nature spéciale, qui se traduisent, sur \inserted{chacun} des torseurs $R_a$ $(a\in A)$ \unsure{connus} sous
$\mathbb Z_a/r(a)\deleted{\unsure{\mathbb Z\mathbin{\wedge}\omega(a)}}$
$(\mathbb Z_a=\mathbb Z\mathbin{\wedge}_{\{\pm1\}}\omega(a))$,
par une opération de contraction des groupes d'opérateurs de $\mathbb Z_a/r(a)$ à $\mathbb Z_a/r'(a)$, où $r'(a)\mid r(a)$
(opération de \uline{gommage partiel} \inserted{\uline{uniforme}} \unsure{en} $a$)
ou d'extension des groupes d'opérateurs de $\mathbb Z_a/r(a)$ à $\mathbb Z_a/\unsure{r''(a)|r'(a)}$, où $\unsure{r''(a)|r'(a)}$ \inserted{$=d(a)r(a)$} est un multiple de $r(a)$
(opération de \uline{surmarquage uniforme en $a$}, \deleted{)} de multiplicateur $d(a)$).
On se rend compte \inserted{ici} que le gommage \uuline{partiel} est en fait un renforcement de structure, pas un affaiblissement, et inversement pour le surmarquage uniforme --- qui, au niveau combinatoire, sera défini sans donnée de structure supplémentaire (sauf bien sûr les multiplicateurs $d(a)$).
\mspage{56}
